\section{Two Layers of Vendor, and Why the Distinction Matters}

\textbf{Why this matters:} A tooling vendor's product inherits the underlying model's limitations and can only ever add mitigations on top of them. Confusing the two hides who is actually capable of fixing a given failure --- and who is not.

\textbf{Practical implication:} When something goes wrong, ask first whether the vendor's own product architecture caused it, or whether it's a limitation of the underlying model that no amount of that vendor's engineering will fix.

Before anything else can be reasoned about clearly, one terminological confusion needs resolving, because it silently collapses two different claims into one word.

\textbf{Model vendors} --- Anthropic, OpenAI, Google, DeepSeek, and similar organizations --- build large language models and their native harnesses: the decoding stack, safety training, tool-use format, default scaffolding (reasoning traces, and so on).

\textbf{Tooling vendors} build products around one or more models, wrapping them in an application-specific harness --- retrieval, verification logic, control flow, an action layer. A contract-review product or a customer-service agent is a tooling-vendor product in the large majority of cases.

The distinction is useful, but the boundary is not as clean in practice as a two-tier description suggests. Some tooling vendors call a model vendor's API and do essentially nothing else; others self-host an open-weight model; others fine-tune or distill a model for their specific task; and a smaller number train a specialist model largely from scratch, sometimes combining several of these approaches within one product. The defensible general statement is that many AI products depend substantially on a third-party foundation model, although the form and depth of that dependency varies considerably --- and the appropriate evaluation method shifts accordingly. A pure API wrapper's quality is dominated by model choice and harness design; a product built on a heavily fine-tuned or distilled model needs to be evaluated on that fine-tuning's own merits, separately from the base model's.

The distinction matters for deciding \textbf{who is responsible for which failure and which claims need which kind of check.} A claim about raw model capability, safety testing, or training data is a \textbf{model-vendor claim}, best tested by independent benchmarking or red-teaming. A claim about how a specific product behaves --- "we verify every citation," "the agent runs in an isolated sandbox" --- is a \textbf{tooling-vendor claim}, best tested by a direct teardown of that product's actual architecture.

A tooling vendor cannot change the \emph{unmodified base model's weights} --- the raw model, called with no scaffolding, has whatever failure profile it has. But a tooling vendor's system-level choices --- retrieval, constrained decoding, deterministic tool calls, targeted fine-tuning, routing between models, restricting what actions the system is allowed to take --- can eliminate or substantially reduce the \emph{deployed system's} failure rate on a defined task, even without changing the base model at all. The useful distinction is between \textbf{limitations of the unchanged base weights} (a model-vendor-layer property) and \textbf{the failure rate of the deployed system as actually built} (a tooling-vendor-layer property, which can differ enormously from the base model's raw behavior). Conflating the two erases exactly the leverage a tooling vendor has to offer.

Product performance is jointly determined by the model, the data it's given, the harness built around it, the controls in place, and the operating context --- and which of these dominates varies by task. For some tasks, model choice swamps everything else; for others, orchestration and retrieval quality matter more than which frontier model sits underneath. This is the core of why an AI tool is not "an AI": it is software that wraps a model, and the wrapping is a real, evaluable, often-decisive engineering artifact in its own right --- not because the wrapping is \emph{always} the dominant factor, but because treating "the model" and "the product" as the same object erases exactly the part of the system a tooling vendor actually built and is accountable for.

Once this is visible, "we've built an autonomous agent" often resolves into something much more mundane: "we've built a loop that calls a model and does something with the results." That's not nothing --- the loop can be excellent or terrible engineering --- but it's an ordinary artifact that can be evaluated on its own merits.

The model itself is a replaceable dependency in principle but rarely operationally or legally fungible in a given deployment. Which model vendor and model a tooling vendor has chosen affects accuracy, context capacity, data terms and residency, version stability under silent updates, indemnities, and concentration risk. Evaluate the tooling vendor's wrapping and the model-vendor supply chain behind it as \textbf{two separate questions} --- a tooling vendor can have excellent verification architecture sitting on a mediocre or unstable model choice, or the reverse.

A fine-tuning claim deserves a precise treatment. The defensible proposition is narrower than blanket skepticism: ordinary fine-tuning is often an unreliable way to inject large amounts of \emph{new} factual knowledge into a model, and retrieval frequently performs better for that specific purpose. But this depends heavily on what kind of fine-tuning is being discussed --- supervised task adaptation, continued pretraining, preference tuning, and deliberate factual-knowledge injection are different techniques with different track records, and fine-tuning \emph{can} improve both factuality and task accuracy when the training objective and data are specifically designed to do that. A tooling vendor's claim ("we trained it on 500,000 filings") should be evaluated on the evidentiary question --- does this measurably improve accuracy on real tasks, or only confidence --- not dismissed by a blanket rule about what fine-tuning generally does.
